\origpage{243}{GRAPH MINORS. XII}
\sectionhead{2. The Trace of a Closed Curve}

Let $\Gamma$ be a 2-cell drawing in a connected surface $\Sigma$, and let $K$ be a radial drawing of $\Gamma$. Let $\phi \colon S^1 \to \Sigma$ be continuous. ($S^1$ is the unit circle $\{(x, y) \colon x^2 + y^2 = 1\}$.) If $\{x \in S^1 \colon \phi(x) \in U(\Gamma)\}$ is finite, we call its cardinality the \textit{length} of $\phi$ (with respect to $\Gamma$). We denote $\{\phi(x) \colon x \in S^1\}$ by $|\phi|$. Now let $\phi$ have finite length and let $|\phi|$ be $\Gamma$-normal. We define a \textit{trace} of $\phi$ in $K$ as follows. If $\phi$ has length 0, its only trace is the zero-edge walk with one vertex $v$, where $v \in V(K)$ is the vertex of $K$ in the region of $\Gamma$ including $|\phi|$. If $\phi$ has length $n > 0$, we proceed as follows. Let
\begin{equation*}
S^1 = I_1 \cup \dots \cup I_n \cup \{x_1, \dots, x_n\},
\end{equation*}
where $x_1, \dots, x_n \in S^1$ are distinct, $\{x_1, \dots, x_n\} = \{x \in S^1 \colon \phi(x) \in V(\Gamma)\}$, $I_1, \dots, I_n$ are mutually disjoint, arc-wise connected open subsets of $S^1 - \{x_1, \dots, x_n\}$, and for $1 \le i \le n$, $\bar{I}_i - I_i = \{x_i, x_{i+1}\}$ (where $x_{n+1}$ means $x_1$). Let $W$ be the closed walk $v_0, e_1, v_1, e_2, \dots, e_{2n}, v_{2n} = v_0$ in $K$, defined by

\hspace*{12pt}(i) For $1 \le i \le n$, $v_{2i} = \phi(x_i)$, and $v_0 = v_{2n}$.

\hspace*{12pt}(ii) For $1 \le i \le n$, $v_{2i-1}$ is the vertex of $K$ in $r_{2i-1}$, where $r_{2i-1}$ is the region of $\Gamma$ including $\{\phi(x) \colon x \in \delm{I_i}\fix{I_{i-1}}\}$\fixnote{fix:2-trace}{Here $v_{2i-1}$ lies between $v_{2i-2} = \phi(x_{i-1})$ and $v_{2i} = \phi(x_i)$, so its region includes $\phi(I_{i-1})$, where $I_0$ means $I_n$. Clause (iii) uses $I_{i-1}$ in the same way. With this indexing, (2.1) holds.}.

\hspace*{12pt}(iii) For $1 \le i \le n$, $e_{2i}$ is the edge of $K$ joining $v_{2i-1}$ and $v_{2i}$ which enters $v_{2i}$ ``in the same way as $\phi$.'' More precisely, choose $h \in I_{i-1}$; then for any $e \in E(K)$ joining $v_{2i-1}$ and $v_{2i}$, we can construct a continuous $\phi_e$ consisting of the restriction of $\phi$ to the portion of $I_{i-1}$ between $h$ and $x_i$, concatenated with a curve tracing $e$ from $v_{2i}$ to $v_{2i-1}$. There is a unique edge $e$ such that $\phi_e$ is homotopic in $r_{2i-1} \cup \{v_{2i}\}$ to a curve in $r_{2i-1}$, and we define $e_{2i}$ to be this edge $e$. (We \del{omit}\fix{give below} the proof that $e$ exists and is unique and does not depend on $h$.)

\hspace*{12pt}(iv) Similarly, for $1 \le i \le n$, $e_{2i-1}$ is the edge of $K$ joining $v_{2i-2}$ and $v_{2i-1}$ which leaves $v_{2i-2}$ in the same way as $\phi$.

Such a walk $W$ is said to be a \textit{trace} of $\phi$.

\fix{We prove the assertion omitted in (iii). The assertion used in (iv) follows from it by running $\phi$ backwards. Let $v = v_{2i}$ and $r = r_{2i-1}$, and suppose that $\Gamma$ has an edge. (If $\Gamma$ has no edge, then by condition (iv) on page 242, $K$ has no edge, and so a curve of positive length has no trace.) Then $v$ is an end of an edge of $\Gamma$. For otherwise a small circle around $v$ bounds a closed disc in the region of $\Gamma$ around $v$, since this region is an open disc. These two closed discs have the same boundary and are different, since only the small one contains $v$, and so their union is open and closed in $\Sigma$. Since $\Sigma$ is connected, they cover $\Sigma$, and so $U(\Gamma) = \{v\}$, a contradiction.

For an edge $e$ of $K$ joining $v_{2i-1}$ and $v$, and for $h \in I_{i-1}$, let us say that $e$ \textit{enters $v$ like $\phi$ from $h$} if $\phi_e$, formed with this $h$, is homotopic in $r \cup \{v\}$ to a curve in $r$. Since $\phi$ maps $I_{i-1}$ into $r$, adding to $\phi_e$ or removing from it an initial part of $\phi$ in $r$ does not change this property. So if $e$ enters $v$ like $\phi$ from one $h$, then it does so from every $h \in I_{i-1}$.

First we show that some edge $e$ enters $v$ like $\phi$ from some $h$. Suppose first that at least two edges of $\Gamma$ are incident with $v$, a loop counted twice. Then the edges of $\Gamma$ at $v$ divide a small closed disc around $v$ into closed sectors, each a closed disc with $v$ in its boundary, and the points of the small disc not in $U(\Gamma)$ are interior points of the sectors. Since $\phi$ maps the part of $I_{i-1}$ near $x_i$ into $r$, which misses $U(\Gamma)$, and this part is connected, $\phi$ maps it into the interior of one sector $Q$. By condition (iv) on page 242, the edges of $\Gamma$ and of $K$ at $v$ alternate in the cyclic order around $v$, and so some edge $e$ of $K$ at $v$ starts inside $Q$. Let $P$ be the set formed by $v$ and the interior points of $Q$ near $v$. It is simply connected, since it contracts to $v$, and $P - \{v\}$ is path connected. Also $P - \{v\}$ misses $U(\Gamma)$ and contains points of $\phi(I_{i-1}) \subseteq r$, and so it lies in $r$. Since $e$ is connected and misses $U(\Gamma)$, and its points near $v$ are in $P - \{v\}$, it lies in $r$, and so its other end is $v_{2i-1}$, the vertex of $K$ in $r$. Choose $h \in I_{i-1}$ such that $\phi$ maps the portion of $I_{i-1}$ between $h$ and $x_i$ into $P$, choose a point $z$ of $e$ such that $e$ runs inside $P$ from $v$ to $z$, and join $\phi(h)$ to $z$ by a path in $P - \{v\}$. Since $P$ is simply connected, the restriction of $\phi$ to the portion of $I_{i-1}$ between $h$ and $x_i$, followed by a curve tracing $e$ from $v$ to $z$, is homotopic in $P$ to this path. Hence $\phi_e$ is homotopic in $r \cup \{v\}$ to this path followed by a curve tracing $e$ from $z$ to $v_{2i-1}$, which is a curve in $r$. So $e$ enters $v$ like $\phi$ from $h$.

Now suppose that just one edge $f$ of $\Gamma$ is incident with $v$, and it is not a loop. By condition (iv) on page 242, exactly one edge $e$ of $K$ is incident with $v$. Remove from a small disc around $v$ a short initial part of $f$. The rest of this disc is a simply connected set $P$, and $P - \{v\}$ is path connected and misses $U(\Gamma)$. It meets $r$, since $v$ is in the closure of $r$, and so it lies in $r$. The points of $e$ near $v$ are in $P - \{v\}$, and so, as before, $e$ lies in $r$ and joins $v$ and $v_{2i-1}$. Also $\phi$ maps the part of $I_{i-1}$ near $x_i$ into $P$, since this part is mapped near $v$ and into $r$. So the same argument with this set $P$ shows that $e$ enters $v$ like $\phi$.

Next let $e$, $e'$ be edges of $K$ joining $v_{2i-1}$ and $v$ that enter $v$ like $\phi$, from $h$ and from $h'$ say. By the second paragraph, both enter $v$ like $\phi$ from $h$. So $\phi_e$ and $\phi_{e'}$, formed with $h$, are homotopic in $r \cup \{v\}$ to curves in $r$ with the same ends $\phi(h)$ and $v_{2i-1}$. These two curves are homotopic in $r$, since $r$ is an open disc, and so $\phi_e$ and $\phi_{e'}$ are homotopic in $r \cup \{v\}$. Cancelling their common initial part, we find that the curves tracing $e$ and $e'$ from $v$ to $v_{2i-1}$ are homotopic in $r \cup \{v\}$. If $e \ne e'$, then $\bar{e} \cup \bar{e}'$ is an $O$-arc null-homotopic in $r \cup \{v\}$, and so by [5, Theorem (7.2)] it bounds a disc $\Delta \subseteq r \cup \{v\}$. Since $\Delta - \{v\} \subseteq r$ misses $U(\Gamma)$, no edge of $\Gamma$ at $v$ enters $\Delta$, and so two edges of $K$ are consecutive in the cyclic order around $v$, contrary to condition (iv) on page 242. So $e = e'$, which proves that $e$ is unique and does not depend on $h$.}\fixnote{fix:2-enters-proof-added}{The paper omits this proof, and we give it above.}

The main result of this section is the following.

\stmt
(2.1) \textit{Let $\Gamma$ be a 2-cell drawing in a connected surface $\Sigma$, with a radial drawing $K$. Let $\mathscr{T}$ be a respectful tangle in $\Gamma$ of order $\theta \ge 1$, let $\phi \colon S^1 \to \Sigma$ be continuous, of length $<\theta$, such that $|\phi|$ is $\Gamma$-normal, and let $W$ be a trace of $\phi$. Then $ins(W)$ is defined. Moreover, for any atom $c$ of $\Gamma$, if there is an $O$-arc $F \subseteq |\phi|$ with $c \cap ins(F) \ne \emptyset$, then $c \cap ins(W) \ne \emptyset$.}

\textit{Proof.} $W$ has $2n < 2\theta$ edges, where $\phi$ has length $n$, and so $ins(W)$ is defined. Let $F \subseteq |\phi|$ be an $O$-arc with $c \cap ins(F) \ne \emptyset$; and let $\phi' \colon S^1 \to F$ be a homeomorphism. If $W'$ is a trace of $\phi'$, then $K | W'$ is a subgraph of
